\documentclass[UTF8,10pt,twocolumn]{ctexart}

\usepackage[a4paper,margin=1.8cm,columnsep=0.65cm]{geometry}
\usepackage{times}
\usepackage{graphicx}
\usepackage{booktabs}
\usepackage{amsmath}
\usepackage{amssymb}
\usepackage{hyperref}
\usepackage{xcolor}

\hypersetup{
    colorlinks=true,
    linkcolor=blue!50!black,
    urlcolor=blue!50!black,
    citecolor=blue!50!black,
}

\setlength{\parskip}{0.25em}
\setlength{\parindent}{1.5em}

\newcommand{\method}{CACAF-HGSOAD}
\newcommand{\figone}{../work_dirs/figure/Gemini_Generated_Image_eolgi7eolgi7eolg.png}
\newcommand{\figtwo}{../work_dirs/figure/Gemini_Generated_Image_oeqo50oeqo50oeqo.png}
\newcommand{\figthree}{../work_dirs/paper_figures/figure3_formal_selected5/figure3_qualitative_mm_sam_adapter_selected5.png}
\newcommand{\figfour}{../work_dirs/paper_figures/figure4_final/figure4_combined.png}

\title{\bfseries Condition-Aware Cross-Modal Adaptive Fusion for Drone-View Multimodal Semantic Segmentation}
\author{Anonymous Submission}
\date{}

\begin{document}

\maketitle

\begin{abstract}
无人机场景下的多模态语义分割同时面临两类关键挑战。第一，飞行姿态、光照和天气变化会导致 RGB 与热红外模态的信息质量在不同场景中动态波动，静态融合策略难以稳定适配。第二，俯视视角会显著压缩目标像素面积，使交通灯、行人、摩托车等小目标在解码过程中更容易被淹没。针对这两个问题，本文提出 \method{} 框架。具体而言，本文设计了 CACAF（Condition-Aware Cross-Modal Adaptive Fusion），通过模态质量评估与跨模态增强实现条件感知的自适应融合；同时设计了 HGSOAD（Hierarchical Guided Small-Object-Aware Decoder），利用分层解码、通道门控和辅助监督增强小目标恢复能力。该框架以冻结的 SAM2 Hiera-L 作为 RGB 主干，以全量可训练的 ConvNeXt-Tiny 编码热红外模态，在保留基础模型通用表征能力的同时实现面向任务的高效适配。

在 FMB 测试集上，本文方法取得 \textbf{68.62 mIoU}，相较 MM SAM-Adapter 的 \textbf{66.10} 提升 \textbf{+2.52}。统一训练配置下的模块消融进一步验证了方法有效性：去除 CACAF 性能下降 \textbf{5.76 mIoU}，去除 SAGU 性能下降 \textbf{3.05 mIoU}。损失函数消融表明 \textbf{CE + Dice} 优于 \textbf{CE + OHEM} 和 \textbf{CE + Dice + OHEM}。此外，在严格 9 类协议下，本文方法在 MFNet 上取得 \textbf{54.73 mIoU}，与 EGFNet 接近，并优于 MFNet 与 RTFNet 等经典方法。
\end{abstract}

\section{引言}

多模态语义分割通过融合可见光与辅助模态信息，可以在弱光、遮挡和复杂环境下提升场景理解鲁棒性。对于无人机平台，这一任务更具挑战性。其一，俯视视角显著压缩目标的空间尺度，小目标类别更容易在下采样和解码过程中丢失。其二，无人机飞行高度、拍摄姿态和环境条件持续变化，不同模态的可靠性也随场景而动态变化，固定融合策略往往难以充分利用跨模态互补性。

近年来，基础模型在视觉表征方面展现出强泛化能力。Segment Anything Model（SAM）及其后续版本 SAM 2 提供了强大的通用视觉特征，而 ViT-Adapter、SAM-Adapter 和 MM SAM-Adapter 则表明，冻结大模型主干并引入轻量适配模块，是将基础模型迁移到下游密集预测任务的一条有效路径。然而，对 FMB 这类无人机 RGB-T 语义分割场景而言，现有方案仍存在两点不足。首先，它们通常缺乏对输入级模态质量差异的显式建模，难以根据当前场景动态调节两种模态的贡献。其次，解码端往往缺少针对小目标恢复的专门设计，容易在高层语义融合后进一步损失细粒度结构。

为解决上述问题，本文提出 \method{}。在编码端，本文采用 SAM2 Hiera-L + Bottleneck Adapter 作为 RGB 主干，并使用 ConvNeXt-Tiny 编码热红外模态，构建非对称双流架构。在融合端，本文提出 CACAF，通过模态质量评估和双向跨模态注意力实现条件感知融合。在解码端，本文提出 HGSOAD，通过层级 U-Net 式上采样、SAGU 门控和辅助监督增强小目标恢复能力。

本文主要贡献如下。
\begin{enumerate}
    \item 提出 CACAF，通过模态质量评估和跨模态增强实现条件感知的自适应融合。
    \item 提出 HGSOAD 与 SAGU，针对无人机俯视场景中的小目标恢复问题设计轻量解码增强机制。
    \item 在 FMB 测试集上达到 \textbf{68.62 mIoU}，相较 MM SAM-Adapter 提升 \textbf{+2.52}。
    \item 提供统一口径的模块消融和损失函数消融，明确 CACAF、SAGU 及 CE + Dice 训练目标的独立作用。
\end{enumerate}

\section{相关工作}

\subsection{多模态语义分割}

RGB-X 语义分割通过引入热红外、深度、事件或 LiDAR 等辅助模态，提升复杂环境下的鲁棒性。早期 RGB-T 方法多采用双分支卷积编码器和显式融合解码器，例如 MFNet、RTFNet 与 FuseSeg，其核心思路是在编码或解码阶段融合 RGB 与热红外特征。这类方法奠定了 RGB-T 语义分割的基本范式，但对跨模态交互的建模能力仍然有限。

随后，研究重点逐步转向更强的跨模态对齐和特征交互。代表性方法包括 FEANet、GMNet、DooDLeNet 和 CRM 等。与此同时，Transformer 框架被引入 RGB-X 语义分割后，CMX 通过跨模态 Transformer 进行特征交换，CMNeXt 则进一步提出面向 arbitrary-modal 场景的统一分割框架，并在 DeLiVER、MFNet 等基准上取得较强表现。尽管这些方法持续推动了 RGB-X 分割性能，但大多数工作主要面向车载或通用道路场景，针对无人机俯视场景中的模态质量波动与小目标恢复问题仍缺乏专门设计。

\subsection{基础模型适配与密集预测迁移}

基础模型在下游视觉任务中的迁移已成为研究热点。ViT-Adapter 表明，在冻结大规模预训练视觉 Transformer 的前提下，引入轻量适配器即可显著增强密集预测性能。SAM-Adapter 进一步证明，SAM 在若干原生表现不足的场景中可以通过轻量适配完成有效迁移。沿着这一思路，MM SAM-Adapter 将 SAM 主干扩展到多模态语义分割任务，通过多模态融合模块和侧路适配结构，把辅助模态信息注入 RGB 主干特征。

\subsection{无人机场景与小目标分割}

与地面视角 RGB-T 分割相比，无人机场景至少存在两点额外困难。首先，俯视视角使大量交通参与体和基础设施目标只占据极少像素，小目标类别更容易在多层下采样后被淹没。其次，无人机在全天候与复杂照明条件下工作时，RGB 与热红外模态的可用性波动更明显，样本级质量差异更大。因此，针对无人机场景的模型不仅需要有效融合多模态信息，还需要在解码阶段显式保留和恢复细粒度结构。

\section{方法}

\subsection{整体架构}

给定 RGB 图像 $I^{rgb}$、热红外图像 $I^{thm}$ 以及像素级标注 $Y$，目标是学习映射
\begin{equation}
\mathcal{F}:(I^{rgb}, I^{thm}) \mapsto \hat{Y},
\end{equation}
其中 $\hat{Y}$ 表示语义分割预测。本文框架由四部分构成：冻结的 RGB 主干、可训练的热红外编码器、多尺度 CACAF 融合模块以及 HGSOAD 解码器。

\begin{figure*}[t]
    \centering
    \includegraphics[width=\textwidth]{\figone}
    \caption{Overall architecture of \method{}. The framework adopts an asymmetric dual-encoder design, performs condition-aware cross-modal fusion at four scales, and decodes the fused features with a hierarchical small-object-aware decoder.}
    \label{fig:overview}
\end{figure*}

如图~\ref{fig:overview} 所示，RGB 分支使用冻结的 SAM2 Hiera-L 作为主要语义表征来源，以保留基础模型的大规模预训练知识；热红外分支则采用全量可训练的 ConvNeXt-Tiny，以增强模型对低照度、遮挡和成像退化条件的适应性。两路特征在四个尺度上分别进入 CACAF 进行条件感知融合，随后由 HGSOAD 自顶向下逐级恢复空间细节。

\subsection{RGB 编码器与热红外编码器}

本文使用 SAM2 Hiera-L 作为 RGB 主干，并冻结其原始骨干参数，仅训练插入在各个 Hiera block 前的 Bottleneck Adapter。设输入 token 为 $x$，Adapter 输出为
\begin{equation}
\mathrm{Adapter}(x)=W_2(\mathrm{GELU}(W_1x)),
\end{equation}
则进入原始 block 的特征为 $x+\mathrm{Adapter}(x)$。热红外分支采用 ConvNeXt-Tiny，提取四级多尺度特征。该设计形成“强 RGB 主干 + 可塑辅助分支”的非对称多模态结构。

\subsection{CACAF：条件感知跨模态自适应融合}

对于第 $i$ 个尺度的 RGB 特征 $f_i^{rgb}$ 和热红外特征 $f_i^{thm}$，CACAF 包含四个步骤：
\begin{enumerate}
    \item \textbf{通道对齐}
    \begin{equation}
    \hat{f}_i^{rgb}=\psi_i^{rgb}(f_i^{rgb}), \quad \hat{f}_i^{thm}=\psi_i^{thm}(f_i^{thm}).
    \end{equation}
    \item \textbf{模态质量评估}
    \begin{equation}
    q_i^{rgb}=\mathrm{MLP}(\mathrm{GAP}(\hat{f}_i^{rgb})), \quad
    q_i^{thm}=\mathrm{MLP}(\mathrm{GAP}(\hat{f}_i^{thm})),
    \end{equation}
    并进一步得到
    \begin{equation}
    [w_i^{rgb}, w_i^{thm}] = \mathrm{softmax}([q_i^{rgb}, q_i^{thm}]).
    \end{equation}
    \item \textbf{跨模态增强}：对高分辨率的 $f1$ 层，由于 token 数达到 16K，本文仅保留质量加权而不引入 cross-attention；对 $f2/f3/f4$ 三层，则采用双向多头注意力进行跨模态增强。
    \item \textbf{自适应融合}
    \begin{equation}
    f_i^{fused}=w_i^{rgb}\tilde{f}_i^{rgb}+w_i^{thm}\tilde{f}_i^{thm}.
    \end{equation}
\end{enumerate}

\begin{figure}[t]
    \centering
    \includegraphics[width=\columnwidth]{\figtwo}
    \caption{Detailed design of CACAF. At each scale, modality features are aligned to a shared embedding space, weighted by modal quality assessment, enhanced by bidirectional cross-modal attention when affordable, and finally fused with adaptive weights.}
    \label{fig:cacaf}
\end{figure}

图~\ref{fig:cacaf} 进一步展示了 CACAF 的内部流程。该模块并不是简单地对两路特征做拼接或平均，而是首先通过通道对齐建立统一嵌入空间，再利用轻量 MQA 分支预测模态质量权重，并在计算可承受的尺度上通过双向 cross-attention 实现显式的信息交换。

\subsection{HGSOAD：层级引导小目标感知解码}

HGSOAD 采用 U-Net 风格的自顶向下层级解码结构。首先通过 $1\times1$ 卷积把四个尺度统一到相同解码通道数，再通过逐级上采样和 skip connection 恢复空间分辨率。在每个尺度上，HGSOAD 通过 SAGU 执行轻量通道门控：
\begin{equation}
g_i=\sigma\left(\phi(\mathrm{GAP}(x_i))+\phi(\mathrm{GMP}(x_i))\right),
\end{equation}
\begin{equation}
\hat{x}_i = g_i \odot x_i,
\end{equation}
其中 $\phi$ 表示 $1\times1$ Conv $\rightarrow$ ReLU $\rightarrow$ $1\times1$ Conv。该门控形式更适合强调局部显著响应，对小目标类别尤其有利。

\subsection{训练目标}

本文采用主头与辅助头的 CE + Dice 作为训练目标：
\begin{equation}
\mathcal{L}=\mathcal{L}_{main} + \lambda \left(\mathcal{L}_{aux2}+\mathcal{L}_{aux3}\right),
\end{equation}
其中
\begin{equation}
\mathcal{L}_{main}=\mathcal{L}_{CE}(p,y)+\mathcal{L}_{Dice}(p,y),
\end{equation}
\begin{equation}
\mathcal{L}_{auxk}=\mathcal{L}_{CE}(p_k^{aux},y)+\mathcal{L}_{Dice}(p_k^{aux},y), \quad k\in\{2,3\},
\end{equation}
并设置 $\lambda=0.4$。实验表明，Dice 与 OHEM 同时使用会导致性能下降，因此最终主结果采用 CE + Dice。

\section{实验}

\subsection{数据集与评测协议}

FMB 是一个面向真实复杂环境的 RGB-T 语义分割基准，包含 14 个语义类别和 $1060/160/280$ 的 train/val/test 划分，原始测试分辨率为 $800\times600$。训练阶段使用随机缩放和随机裁剪后的 $512\times512$ 输入；主结果评测则在原始测试分辨率下进行，仅在推理前将图像 pad 到 32 的整数倍，而不对整图强制缩放到 $512\times512$。

\subsection{实现细节}

优化器采用 AdamW，SAM2 Adapter 与 ConvNeXt-Tiny 的学习率均设为 $1\times10^{-4}$，CACAF 与 HGSOAD 的学习率设为 $2\times10^{-4}$，权重衰减为 $0.01$。学习率调度采用 500 iter warmup 加 cosine decay。模型在 BF16 混合精度下训练 60 个 epoch，batch size 为 4。数据增强包括随机缩放、随机裁剪、随机水平翻转、Gaussian blur（$p=0.2$）、photometric distortion 以及 $cat\_max\_ratio=0.75$ 约束。

\subsection{与现有方法对比}

\begin{table}[t]
    \centering
    \caption{Comparison with the direct baseline on FMB test.}
    \label{tab:main}
    \begin{tabular}{lccc}
        \toprule
        Method & mIoU & mAcc & aAcc \\
        \midrule
        MM SAM-Adapter & 66.10 & 72.09 & 93.95 \\
        \textbf{Ours (\method)} & \textbf{68.62} & \textbf{75.75} & 93.27 \\
        \bottomrule
    \end{tabular}
\end{table}

表~\ref{tab:main} 给出了本文方法与直接基线 MM SAM-Adapter 在 FMB 测试集上的对比结果。相较基线，本文实现了 \textbf{+2.52 mIoU} 与 \textbf{+3.66 mAcc} 提升。虽然 aAcc 略低于基线，但本文方法在少样本小目标类别上的提升更显著，说明性能增益主要来自对困难类别和细粒度结构的更好恢复。

\subsection{模块消融与损失函数消融}

\begin{table}[t]
    \centering
    \caption{Module ablation on FMB test.}
    \label{tab:ablation}
    \begin{tabular}{lcc}
        \toprule
        Method & mIoU & $\Delta$ \\
        \midrule
        \textbf{Full (CACAF + SAGU)} & \textbf{68.62} & --- \\
        w/o CACAF & 62.86 & -5.76 \\
        w/o SAGU & 65.57 & -3.05 \\
        w/o CACAF, w/o SAGU & 64.78 & -3.84 \\
        \bottomrule
    \end{tabular}
\end{table}

\begin{table}[t]
    \centering
    \caption{Loss function ablation on FMB test.}
    \label{tab:loss}
    \begin{tabular}{lcc}
        \toprule
        Loss & mIoU & $\Delta$ \\
        \midrule
        \textbf{CE + Dice} & \textbf{68.62} & --- \\
        CE + OHEM & 67.79 & -0.83 \\
        CE + Dice + OHEM & 66.79 & -1.83 \\
        \bottomrule
    \end{tabular}
\end{table}

结果表明，CACAF 是主要性能来源，去除后性能下降最为明显；SAGU 也能在完整框架中带来稳定增益，说明解码端的小目标增强是必要的。损失函数消融进一步表明，在本文网络结构下，Dice 与 OHEM 的组合并未形成有效互补，反而会带来性能退化，因此最终主模型采用 CE + Dice 作为训练目标。

\subsection{跨数据集补充实验}

\begin{table}[t]
    \centering
    \caption{MFNet comparison under the strict 9-class protocol.}
    \label{tab:mfnet}
    \begin{tabular}{lc}
        \toprule
        Method & mIoU \\
        \midrule
        MFNet (2017) & 45.1 \\
        RTFNet (2019) & 53.2 \\
        EGFNet (2022) & 54.8 \\
        GMNet (2021) & 57.3 \\
        \textbf{Ours (\method)} & \textbf{54.73} \\
        \bottomrule
    \end{tabular}
\end{table}

在严格 9 类协议下，本文方法在 MFNet 上取得 \textbf{54.73 mIoU}、\textbf{68.14 mAcc} 和 \textbf{97.90 aAcc}。与经典的 MFNet 和 RTFNet 相比，本文方法表现更优；与 EGFNet 相比结果接近；与 GMNet 相比仍有差距。

\begin{figure*}[t]
    \centering
    \includegraphics[width=\textwidth]{\figthree}
    \caption{Qualitative comparisons on five representative FMB test scenes (\texttt{00779}, \texttt{00048}, \texttt{00121}, \texttt{00712}, and \texttt{01478}). Compared with MM SAM-Adapter, the proposed method yields clearer boundaries and more reliable recovery of small or weakly illuminated targets under challenging RGB-T conditions.}
    \label{fig:qualitative}
\end{figure*}

\begin{figure}[t]
    \centering
    \includegraphics[width=\columnwidth]{\figfour}
    \caption{Visualization of adaptive fusion weights in CACAF. Thermal contributions are concentrated at shallow and intermediate scales, while deeper scales are dominated by the RGB branch.}
    \label{fig:weights}
\end{figure}

\subsection{定性对比与可解释性分析}

图~\ref{fig:qualitative} 给出了 FMB 测试集上的 5 个正式代表性样本（\texttt{00779}、\texttt{00048}、\texttt{00121}、\texttt{00712}、\texttt{01478}），包含 RGB、Thermal、Ground Truth、MM SAM-Adapter 和 Ours 五列对比。可以看到，本文方法在弱光和低对比度场景下能够更稳定地恢复行人、交通灯以及远距离车辆等细粒度目标，同时在道路边界和前景轮廓上保持更清晰的结构。图~\ref{fig:weights} 进一步从权重统计角度分析了 CACAF 的条件感知行为，说明 Thermal 模态在浅层和中层特征上的贡献更明显，而深层语义表示更倾向于由冻结的 SAM2 主干主导。

\section{讨论与结论}

与固定地面视角场景不同，无人机拍摄时的光照、视角和成像距离变化更频繁，RGB 与热红外模态的有效性也更不稳定。CACAF 通过显式质量估计和动态加权，使得融合策略可以随输入样本变化而调整，从而更适合无人机场景。与此同时，FMB 中多类关键目标在俯视视角下仅占据少量像素，单纯依靠高层融合特征并不足以稳定恢复这些细粒度结构。HGSOAD 通过层级解码、SAGU 门控和辅助监督增强浅层细节传播，因此对交通灯、公交、摩托车等类别更有帮助。

总体而言，\method{} 在保持“冻结基础模型主干 + 轻量适配”范式的同时，引入条件感知跨模态融合和小目标导向解码机制，在 FMB 测试集上取得了稳定有效的性能提升。后续工作将进一步补充更完整的效率对比与跨数据集评测，并根据最终投稿目标替换为对应会议模板。

\begin{thebibliography}{99}

\bibitem{sam}
A. Kirillov et al.
\newblock Segment Anything.
\newblock In \emph{ICCV}, 2023.

\bibitem{sam2}
N. Ravi et al.
\newblock SAM 2: Segment Anything in Images and Videos.
\newblock arXiv preprint arXiv:2408.00714, 2024.

\bibitem{vitadapter}
Z. Chen et al.
\newblock Vision Transformer Adapter for Dense Predictions.
\newblock arXiv preprint arXiv:2205.08534, 2022.

\bibitem{samadapter}
T. Chen et al.
\newblock SAM-Adapter: Adapting Segment Anything in Underperformed Scenes.
\newblock arXiv preprint arXiv:2304.09148, 2023.

\bibitem{cmx}
J. Zhang et al.
\newblock CMX: Cross-Modal Fusion for RGB-X Semantic Segmentation with Transformers.
\newblock arXiv preprint arXiv:2203.04838, 2022.

\bibitem{cmnext}
J. Zhang et al.
\newblock Deliving Arbitrary-Modal Semantic Segmentation.
\newblock In \emph{CVPR}, 2023.

\bibitem{mmsam}
F. Curti et al.
\newblock MM SAM-Adapter: Multimodal Semantic Segmentation with SAM.
\newblock arXiv preprint arXiv:2509.10408, 2025.

\bibitem{fmb}
Q. Ha et al.
\newblock MFNet and FMB: Multispectral Semantic Segmentation Benchmarks.
\newblock 2022.

\bibitem{mfnet}
Q. Ha et al.
\newblock MFNet: Toward Real-Time Semantic Segmentation for Autonomous Vehicles with Multispectral Scenes.
\newblock In \emph{IROS}, 2017.

\bibitem{rtfnet}
Y. Sun et al.
\newblock RTFNet: RGB-Thermal Fusion Network for Semantic Segmentation of Urban Scenes.
\newblock \emph{IEEE Robotics and Automation Letters}, 2019.

\end{thebibliography}

\end{document}
